\documentclass[a4paper,12pt]{article}
\usepackage[T2A]{fontenc}
\usepackage[utf8]{inputenc}
\usepackage[russian]{babel}
\usepackage{PTSans}
\renewcommand{\familydefault}{\sfdefault}
\usepackage{amsmath,amssymb,mathtools}
\usepackage{icomma}
\usepackage[a4paper,left=22mm,right=22mm,top=18mm,bottom=18mm,headheight=26pt,headsep=5mm]{geometry}
\usepackage{microtype}
\usepackage{tikz}
\usetikzlibrary{arrows.meta,positioning,shapes.geometric,calc}
\usepackage{booktabs,tabularx}
\usepackage{enumitem}
\usepackage{fancyhdr}
\usepackage{setspace}
\usepackage{xcolor}
\usepackage{hyperref}
\definecolor{accent}{HTML}{087C82}
\definecolor{darkaccent}{HTML}{174F59}
\definecolor{textgray}{HTML}{333D43}
\definecolor{softgray}{HTML}{5C6D75}
\definecolor{lightline}{HTML}{D6E1E1}
\hypersetup{colorlinks=true,linkcolor=darkaccent,urlcolor=darkaccent,pdfauthor={Лёвин Артём Александрович},pdftitle={Равнодействующая сил и второй закон Ньютона — 17.10.26}}
\setstretch{1.14}
\setlength{\parindent}{0pt}
\setlength{\parskip}{2.5pt}
\setlist[itemize]{leftmargin=*,topsep=2pt,itemsep=2pt,parsep=0pt}
\setlist[enumerate]{leftmargin=*,topsep=4pt,itemsep=10pt,parsep=0pt}
\pagestyle{fancy}
\fancyhf{}
\fancyhead[L]{\small\color{softgray}ФИЗИКА \(\cdot\) ОГЭ}
\fancyhead[R]{\small\color{softgray}17.10.26}
\fancyfoot[C]{\small\color{softgray}\thepage}
\renewcommand{\headrulewidth}{0.35pt}
\renewcommand{\headrule}{\hbox to\headwidth{\color{lightline}\leaders\hrule height\headrulewidth\hfill}}
\newcommand{\sectionhead}[1]{\vspace{6pt}{\large\bfseries\color{darkaccent}#1}\par\vspace{2pt}}
\newcommand{\topichead}[1]{\vspace{3pt}{\normalsize\bfseries\color{darkaccent}#1}\par}
\newcommand{\checkitem}{\item[$\square$]}
\newcommand{\emphkey}[1]{\textbf{\color{darkaccent}#1}}
\newcommand{\unit}[1]{\,\text{#1}}
\newcommand{\ruleline}{\par\vspace{2pt}\noindent{\color{lightline}\rule{\linewidth}{0.45pt}}\par\vspace{2pt}}
\tikzset{
  force/.style={-{Stealth[length=3mm]},line width=1.45pt,draw=accent},
  otherforce/.style={-{Stealth[length=3mm]},line width=1.45pt,draw=darkaccent},
  resultant/.style={-{Stealth[length=3.4mm]},line width=1.8pt,draw=black},
  flowstart/.style={ellipse,draw=darkaccent,line width=0.9pt,align=center,minimum width=38mm,minimum height=10mm,fill=accent!5},
  flowaction/.style={rounded corners=2pt,rectangle,draw=darkaccent,line width=0.9pt,align=center,text width=110mm,minimum height=13mm,inner sep=5pt},
  flowdecision/.style={diamond,aspect=2.55,draw=darkaccent,line width=0.9pt,align=center,text width=53mm,inner sep=1.5pt},
  flowarr/.style={-{Stealth[length=2.8mm]},line width=1pt,draw=textgray},
}
\begin{document}
\thispagestyle{empty}
\begin{center}
\vspace*{35mm}
{\small\bfseries\color{accent} ФИЗИКА \(\cdot\) ПОДГОТОВКА К ОГЭ}\par
\vspace{16mm}
{\Huge\bfseries\color{darkaccent}Чек-лист}\par
\vspace{9mm}
{\LARGE\bfseries Равнодействующая сил}\par
\vspace{3mm}
{\Large Второй закон Ньютона и направление ускорения}\par
\vspace{12mm}
{\normalsize\color{textgray}Сложение векторов, силы в механике, проекции}\par
{\normalsize\color{textgray}Повторение кинематики и движения по окружности}\par
\vspace{19mm}
{\large 17 октября 2026 года}\par
\vfill
{\normalsize\color{textgray}Лёвин Артём Александрович}\par
{\normalsize\color{textgray}эксклюзивно для Володи Хачатуряна}\par
\vspace{22mm}
\begin{tikzpicture}[x=1cm,y=1cm]
\draw[accent!70,line width=1.15pt] (-5,0) .. controls (-2.7,0.9) and (-1.3,-0.85) .. (1.0,0) .. controls (2.7,0.6) and (4.1,0.53) .. (5,0.12);
\end{tikzpicture}
\end{center}
\clearpage

{\LARGE\bfseries\color{darkaccent}Основной чек-лист}\par
\ruleline
\sectionhead{1. Кинематика: что вспомнили на занятии}
\begin{itemize}
\checkitem При постоянном ускорении вдоль оси \(Ox\):
\[
\boxed{x=x_0+v_{0x}t+\frac{a_xt^2}{2}},\qquad
\boxed{v_x=v_{0x}+a_xt}.
\]
По записи \(x=-1+4t+5t^2\) (в СИ) определяем: \(x_0=-1\unit{м}\), \(v_{0x}=4\unit{м/с}\), \(a_x=10\unit{м/с\textsuperscript{2}}\). Обязательно учитываем коэффициент \(\tfrac12\) перед ускорением.
\checkitem Для перемещения \(\Delta x=x-x_0\), если ускорение постоянно:
\[
\Delta x=v_{0x}t+\frac{a_xt^2}{2}
=\frac{v_{0x}+v_x}{2}\,t,
\qquad v_x^2-v_{0x}^2=2a_x\Delta x.
\]
Последнее равенство не требует знания времени. Формулу
\(\Delta x=(v_x^2-v_{0x}^2)/(2a_x)\) можно применять при \(a_x\ne0\).
\checkitem Различаем \emphkey{«за первые \(n\) секунд»} и \emphkey{«за \(n\)-ю секунду»}. За \(n\)-ю секунду:
\[
\Delta x_n=x(n\unit{с})-x((n-1)\unit{с}).
\]
Если \(v_0\) задана в м/с, \(a\) в м/с\(^2\), то численное значение перемещения за \(n\)-ю секунду (в метрах):
\[
\boxed{\Delta x_n=v_0+\frac{a(2n-1)}{2}}.
\]
Это выражение для интервала длительностью \(1\unit{с}\). Назвать \(\Delta x_n\) \emph{путём} можно, если на данном интервале нет изменения направления движения.
\end{itemize}
\topichead{Мини-пример}
При \(v_0=2\unit{м/с}\), \(a=2\unit{м/с\textsuperscript{2}}\) за первые \(4\unit{с}\) тело проходит \(s(4)=2\cdot4+2\cdot4^2/2=24\unit{м}\). За четвёртую секунду: \(s(4)-s(3)=24-15=9\unit{м}\).

\sectionhead{2. Равномерное движение по окружности}
\begin{itemize}
\checkitem Период \(T=t/N\) --- \emph{время одного оборота} (с). Частота \(\nu=N/t\) --- \emph{число оборотов за секунду} (Гц). Связь:
\[\boxed{T=\frac1\nu}.\]
\checkitem Линейная скорость \(v=2\pi R/T=2\pi R\nu\) (м/с); угловая скорость \(\omega=2\pi/T=2\pi\nu\) (рад/с).
\checkitem Связь скоростей и модуль центростремительного ускорения:
\[
\boxed{v=\omega R},\qquad
\boxed{a_{\text{цс}}=\frac{v^2}{R}=\omega^2R}.
\]
Центростремительное ускорение направлено \emphkey{к центру окружности}. Радиус \(R>0\).
\end{itemize}
\clearpage

{\LARGE\bfseries\color{darkaccent}Силы и их сложение}\par
\ruleline
\sectionhead{3. Равнодействующая и второй закон Ньютона}
\begin{itemize}
\checkitem \emphkey{Равнодействующая} --- \textbf{векторная сумма всех сил}, приложенных к выбранному телу:
\[\boxed{\vec F_{\text{рез}}=\sum_i\vec F_i}.\]
Силы складываем с учётом их \emph{направлений}; обычная сумма модулей подходит не всегда.
\checkitem Второй закон Ньютона в инерциальной системе отсчёта:
\[
\boxed{\sum_i\vec F_i=m\vec a},\qquad
\boxed{\sum_i F_{iy}=ma_y}.
\]
Направления \(\vec F_{\text{рез}}\) и \(\vec a\) совпадают при \(F_{\text{рез}}\ne0\). \emphkey{Направление скорости может быть другим.}
\checkitem По известным \(m,v_0,v,t\) при постоянном ускорении:
\[a_x=\frac{v_x-v_{0x}}{t},\qquad F_{\text{рез},x}=ma_x.\]
Минус в результате говорит о направлении проекции. Если просят \emph{модуль} силы, ответ неотрицателен.
\end{itemize}
\topichead{Пример: торможение молотка}
Молоток массой \(0{,}8\unit{кг}\) за \(0{,}2\unit{с}\) уменьшил скорость с \(5\unit{м/с}\) до нуля. Ось направлена по начальной скорости:
\[
a_x=\frac{0-5}{0{,}2}=-25\unit{м/с\textsuperscript{2}},\quad
F_{\text{рез},x}=0{,}8\cdot(-25)=-20\unit{Н}.
\]
\emphkey{Модуль средней равнодействующей --- \(20\unit{Н}\).} При пренебрежении остальными силами во время удара это приблизительно модуль средней контактной силы.

\sectionhead{4. Векторы на одной прямой}
\begin{itemize}
\checkitem Если силы направлены в одну сторону, их модули \emphkey{складываются}.
\checkitem Если силы направлены в противоположные стороны, модуль равнодействующей равен \emphkey{разности большего и меньшего модулей}, направление --- к большей силе.
\end{itemize}
\topichead{Пример}
\(F_1=10\unit{Н}\) вниз, \(F_2=20\unit{Н}\) вверх. Равнодействующая \(F_{\text{рез}}=20-10=10\unit{Н}\) \emph{вверх}.

\clearpage

{\LARGE\bfseries\color{darkaccent}Векторное сложение на рисунке}\par
\ruleline
\sectionhead{5. Перпендикулярные силы}
\begin{itemize}
\checkitem Переносим второй вектор параллельно себе так, чтобы его начало совпало с концом первого. Соединяем начало первого с концом второго.
\checkitem Если два вектора взаимно перпендикулярны, модуль их суммы:
\[\boxed{F_{\text{рез}}=\sqrt{F_1^2+F_2^2}}.\]
\end{itemize}
\vspace{3pt}
\sectionhead{6. Пример с тремя силами}
На тело действуют \(\vec F_1\) вверх (\(4\unit{Н}\)), \(\vec F_2\) влево (\(2\unit{Н}\)) и \(\vec F_3\) вправо (\(3\unit{Н}\)). На рисунке масштаб \(1\unit{Н}=0{,}65\unit{см}\).
\begin{center}
\begin{tikzpicture}[x=0.65cm,y=0.65cm]
\draw[step=1,lightline!70,thin] (-3.0,-0.5) grid (4.8,5.2);
\coordinate (O) at (0,0);
\fill[darkaccent] (O) circle (2.3pt);
\draw[force] (O)--(0,4) node[above,black,font=\small] {$\vec F_1=4\unit{Н}$};
\draw[otherforce] (O)--(-2,0) node[below,black,font=\small] {$\vec F_2=2\unit{Н}$};
\draw[force] (O)--(3,0) node[below,black,font=\small] {$\vec F_3=3\unit{Н}$};
\draw[otherforce,dashed] (0,4)--(3,4);
\draw[resultant] (0,0)--(3,4) node[pos=0.64,right=3pt,font=\small,fill=white,inner sep=1pt] {$\vec F_1+\vec F_3$};
\end{tikzpicture}
\end{center}
\vspace{-2pt}
\begin{itemize}
\checkitem \(\vec F_1+\vec F_3\): взаимно перпендикулярны, поэтому \(F_{13}=\sqrt{4^2+3^2}=5\unit{Н}\), направление --- вправо и вверх.
\checkitem \(\vec F_1+\vec F_2\): \(F_{12}=\sqrt{4^2+2^2}=2\sqrt5\unit{Н}\), направление --- влево и вверх.
\checkitem \(\vec F_2+\vec F_3\): противоположные, поэтому \(F_{23}=3-2=1\unit{Н}\) вправо.
\checkitem \emphkey{Все три силы}: по горизонтали остаётся \(1\unit{Н}\) вправо, по вертикали \(4\unit{Н}\) вверх. Следовательно,
\[F_{\text{рез}}=\sqrt{1^2+4^2}=\sqrt{17}\unit{Н}.\]
\end{itemize}
\ruleline
\topichead{Важно различать}
\(\vec F_1+\vec F_2\) --- вектор; \(\lvert\vec F_1+\vec F_2\rvert\) --- его модуль. При вычислениях через теорему Пифагора квадрат стоит над \emph{модулем результирующей силы}: \(F_{12}^2=F_1^2+F_2^2\).
\clearpage

{\LARGE\bfseries\color{darkaccent}Направления сил и проекции}\par
\ruleline
\sectionhead{7. Какие силы могут действовать на тело}
\begin{itemize}
\checkitem \emphkey{Сила тяжести} \(\vec F_{\text{тяж}}=m\vec g\): направлена вертикально вниз; вблизи поверхности Земли \(F_{\text{тяж}}\approx mg\).
\checkitem \emphkey{Реакция опоры} \(\vec N\): направлена перпендикулярно поверхности опоры \emph{от неё к телу}.
\checkitem \emphkey{Сила натяжения нити} \(\vec T\): направлена вдоль натянутой нити \emph{от тела}.
\checkitem \emphkey{Сила трения скольжения}: направлена против относительного скольжения соприкасающихся поверхностей. На неподвижном контакте трение покоя может иметь другое направление, поэтому правило «против скорости тела» универсальным не является.
\checkitem \emphkey{Сила тяги} --- если имеется, направлена в сторону действия двигателя или тягового устройства на тело.
\end{itemize}
\sectionhead{8. Почему важно выбрать ось}
Для оси \(Oy\), направленной вверх, знаки проекций определяются \emph{по направлениям сил и ускорения}, а не по желаемому ответу.

\topichead{Человек в лифте, стоящий на полу}
На человека действуют \(\vec N\) вверх и \(m\vec g\) вниз. В системе отсчёта Земли (приближённо инерциальной):
\[
\boxed{N-mg=ma_y}.
\]
\begin{center}
\begin{tikzpicture}[x=1cm,y=1cm]
\begin{scope}[shift={(-3.25,0)}]
\draw[lightline,line width=1pt] (-1,-0.55)--(1,-0.55);
\draw[fill=accent!7,draw=darkaccent,rounded corners=2pt] (-0.5,-0.52) rectangle (0.5,0.35);
\draw[force] (0,0.35)--(0,1.75) node[above,black] {$\vec N$};
\draw[otherforce] (0,-0.51)--(0,-1.65) node[below,black] {$m\vec g$};
\node[below=24mm,font=\small,align=center] at (0,-0.9) {Ускорение вверх:\\$N=m(g+a)$};
\end{scope}
\begin{scope}[shift={(3.25,0)}]
\draw[lightline,line width=1pt] (-1,-0.55)--(1,-0.55);
\draw[fill=accent!7,draw=darkaccent,rounded corners=2pt] (-0.5,-0.52) rectangle (0.5,0.35);
\draw[force] (0,0.35)--(0,1.38) node[above,black] {$\vec N$};
\draw[otherforce] (0,-0.51)--(0,-2.03) node[below,black] {$m\vec g$};
\node[below=24mm,font=\small,align=center] at (0,-0.9) {Ускорение вниз:\\$N=m(g-a)$};
\end{scope}
\end{tikzpicture}
\end{center}
\vspace{7pt}
В обеих формулах \(a>0\) --- \emph{модуль} ускорения. Для движения вниз при ускорении \(a<g\) реакция опоры остаётся положительной. Если \(a=0\), то \(N=mg\) и лифт может покоиться или двигаться равномерно.
\clearpage

{\LARGE\bfseries\color{darkaccent}Алгоритм: второй закон Ньютона}\par
\ruleline
\vspace{15mm}
\begin{center}
\begin{tikzpicture}[node distance=13mm]
\node[flowstart] (s) {Начало};
\node[flowaction,below=of s] (a) {Выберите тело и изобразите\\все силы, действующие на него};
\node[flowdecision,below=13mm of a] (b) {Все силы\\учтены?};
\node[flowaction,below=17mm of b] (c) {Выберите ось и укажите\\направление ускорения};
\node[flowaction,below=of c] (d) {Запишите проекции: \\ $\sum F_y=ma_y$ (или $\sum F_x=ma_x$)};
\node[flowaction,below=of d] (e) {Подставьте данные в СИ, решите уравнение;\\проверьте знак, направление и единицы};
\node[flowstart,below=of e] (f) {Ответ};
\draw[flowarr] (s)--(a);
\draw[flowarr] (a)--(b);
\draw[flowarr] (b)--node[right,font=\small,fill=white] {да} (c);
\draw[flowarr] (c)--(d);
\draw[flowarr] (d)--(e);
\draw[flowarr] (e)--(f);
\draw[flowarr] (b.west)--++(-26mm,0) node[pos=0.45,above,font=\small] {нет} |- (a.west);
\end{tikzpicture}
\end{center}
\clearpage

{\LARGE\bfseries\color{darkaccent}Тренировочный блок}\par
{\small\color{softgray}Путь А \(\cdot\) 10 заданий \(\cdot\) от базовых к составным}\par
\ruleline
Записывайте формулы, подстановку и ответ с единицами. Для сил указывайте направление, если этого требует условие.
\begin{enumerate}[label=\textbf{\arabic*.},itemsep=12pt]
\item Координата тела в СИ задаётся формулой \(x=-2+3t+4t^2\). Определите начальную координату, проекцию начальной скорости и проекцию ускорения.
\item Тело движется прямолинейно без смены направления с начальной скоростью \(2\unit{м/с}\) и постоянным ускорением \(2\unit{м/с\textsuperscript{2}}\). Найдите путь за первые четыре секунды и путь \emph{за четвёртую секунду}.
\item Точка совершает \(30\) равномерных оборотов по окружности радиуса \(0{,}40\unit{м}\) за \(15\unit{с}\). Найдите период, частоту, угловую и линейную скорости. Ответ для скоростей можно оставить через \(\pi\).
\item Молоток массой \(600\unit{г}\) перед ударом двигался со скоростью \(6\unit{м/с}\) и остановился за \(0{,}15\unit{с}\). Считая ускорение постоянным, найдите его проекцию (ось вдоль начальной скорости) и модуль средней равнодействующей силы за время удара.
\item На тело действуют две горизонтальные силы: \(18\unit{Н}\) вправо и \(11\unit{Н}\) влево. Определите модуль и направление их равнодействующей.
\end{enumerate}
\clearpage
{\LARGE\bfseries\color{darkaccent}Тренировочный блок \(\cdot\) продолжение}\par
\ruleline
\begin{enumerate}[label=\textbf{\arabic*.},start=6,itemsep=17pt]
\item Две силы, приложенные к одной точке, взаимно перпендикулярны: \(6\unit{Н}\) вверх и \(8\unit{Н}\) вправо. Найдите модуль равнодействующей и укажите её направление.
\item На тело действуют три силы: \(9\unit{Н}\) вверх, \(12\unit{Н}\) влево и \(7\unit{Н}\) вправо. Найдите модуль равнодействующей и определите, в какую сторону она направлена (по горизонтали и вертикали).
\item На тело массой \(3\unit{кг}\) действуют силы \(21\unit{Н}\) вверх и \(9\unit{Н}\) вниз. Найдите равнодействующую и ускорение тела. Укажите их направления.
\item На брусок массой \(5\unit{кг}\), движущийся по горизонтальной плоскости вправо, действуют сила тяги \(22\unit{Н}\) вправо и сила трения скольжения \(7\unit{Н}\) влево. По вертикали сила тяжести уравновешена реакцией опоры. Найдите ускорение бруска.
\item Человек массой \(50\unit{кг}\) стоит на полу лифта. Определите модуль реакции опоры, если лифт движется: 
\begin{enumerate}[label=\alph*),leftmargin=8mm,itemsep=3pt,topsep=3pt]
\item с ускорением \(2\unit{м/с\textsuperscript{2}}\) вверх;
\item с ускорением \(2\unit{м/с\textsuperscript{2}}\) вниз.
\end{enumerate}
Примите \(g=10\unit{м/с\textsuperscript{2}}\). В обоих случаях выберите вертикальную ось вверх и запишите второй закон Ньютона в проекциях.
\end{enumerate}
\clearpage

{\LARGE\bfseries\color{darkaccent}Ответы для самопроверки}\par
\ruleline
\vspace{4mm}
\renewcommand{\arraystretch}{1.62}
\begin{tabularx}{\linewidth}{@{}p{15mm}X@{}}
\toprule
\textbf{№} & \textbf{Ответ} \\
\midrule
1 & \(x_0=-2\unit{м}\), \(v_{0x}=3\unit{м/с}\), \(a_x=8\unit{м/с\textsuperscript{2}}\). \\
2 & За первые \(4\unit{с}\): \(24\unit{м}\); за четвёртую секунду: \(9\unit{м}\). \\
3 & \(T=0{,}5\unit{с}\), \(\nu=2\unit{Гц}\), \(\omega=4\pi\unit{рад/с}\), \(v=1{,}6\pi\unit{м/с}\). \\
4 & \(a_x=-40\unit{м/с\textsuperscript{2}}\), \(F_{\text{рез}}=24\unit{Н}\). \\
5 & \(7\unit{Н}\), вправо. \\
6 & \(10\unit{Н}\), вправо и вверх. \\
7 & \(\sqrt{106}\unit{Н}\approx10{,}3\unit{Н}\), влево и вверх. \\
8 & \(F_{\text{рез}}=12\unit{Н}\) вверх, \(a=4\unit{м/с\textsuperscript{2}}\) вверх. \\
9 & \(a=3\unit{м/с\textsuperscript{2}}\), вправо. \\
10 & а) \(N=600\unit{Н}\), \(N-mg=ma\);\\
   & б) \(N=400\unit{Н}\), \(N-mg=-ma\). \\
\bottomrule
\end{tabularx}
\vspace{11mm}
\emphkey{Три короткие самопроверки}
\begin{itemize}
\item Складывали ли Вы силы \emph{векторно}, учитывая направления?
\item Совпадает ли найденное направление ускорения с направлением равнодействующей?
\item Отвечает ли формула именно на вопрос задачи: «за \(n\) секунд» или «за \(n\)-ю секунду»?
\end{itemize}
\vfill
{\small\color{softgray}Лёвин Артём Александрович \(\cdot\) исключительно для Володи Хачатуряна}
\end{document}
